
\documentclass[10pt]{article}
\usepackage[utf8]{inputenc}
\usepackage[T1]{fontenc}
\usepackage{amsmath, amssymb, xcolor, geometry, enumitem, multicol, tcolorbox}

\definecolor{mainblue}{RGB}{0, 102, 204}
\geometry{a4paper, margin=1.2cm, top=1cm, bottom=3cm, footskip=1.2cm}

\begin{document}
\enlargethispage{2.5cm}

% Modern Header
\begin{tcolorbox}[colback=mainblue!5, colframe=mainblue, sharp corners, boxrule=0.5pt]
    \small \textbf{Unit:} undefined \hfill \textbf{Name:} \rule{3cm}{0.4pt} \\
    \small \textbf{Topic:} Variables and Expressions / Order of Operations \hfill \textbf{Date:} \rule{2cm}{0.4pt} \\
    \centering \textbf{\large Homework 2}
\end{tcolorbox}

\vspace{0.2cm}

% MCQ Section
\begin{multicols}{2}
\begin{enumerate}[label=\textbf{Q\arabic*.}, leftmargin=*, itemsep=6pt, parsep=0pt]

  \item \small What is the variable in the expression $2x + 3$?
  \begin{enumerate}[label=(\Alph*), leftmargin=0.6cm, nosep, itemsep=2pt]
    \item 2
    \item 3
    \item x
    \item 2x
    \item None of the above
  \end{enumerate}

  \item \small Identify the operator in the expression $4 	imes 5$?
  \begin{enumerate}[label=(\Alph*), leftmargin=0.6cm, nosep, itemsep=2pt]
    \item +
    \item -
    \item $	imes$
    \item $div$
    \item =
  \end{enumerate}

  \item \small What is the coefficient of $x$ in $4x$?
  \begin{enumerate}[label=(\Alph*), leftmargin=0.6cm, nosep, itemsep=2pt]
    \item 1
    \item 2
    \item 3
    \item 4
    \item None of the above
  \end{enumerate}

  \item \small What does the term 'expression' mean in math?
  \begin{enumerate}[label=(\Alph*), leftmargin=0.6cm, nosep, itemsep=2pt]
    \item A statement that is always true
    \item A statement that is always false
    \item A group of numbers, variables, and operators
    \item A single number or value
    \item A mathematical equation
  \end{enumerate}

  \item \small What is the constant term in $2x + 5$?
  \begin{enumerate}[label=(\Alph*), leftmargin=0.6cm, nosep, itemsep=2pt]
    \item 2
    \item 5
    \item x
    \item 2x
    \item None of the above
  \end{enumerate}

  \item \small Identify the operand in $3 + 2$.
  \begin{enumerate}[label=(\Alph*), leftmargin=0.6cm, nosep, itemsep=2pt]
    \item +
    \item 3
    \item 2
    \item Both 3 and 2
    \item None of the above
  \end{enumerate}

  \item \small What does the term 'variable' mean in math?
  \begin{enumerate}[label=(\Alph*), leftmargin=0.6cm, nosep, itemsep=2pt]
    \item A symbol that represents a fixed value
    \item A symbol that represents a value that can change
    \item A mathematical operation
    \item A numerical value
    \item A constant term
  \end{enumerate}

  \item \small What is $2 cdot 3$ an example of?
  \begin{enumerate}[label=(\Alph*), leftmargin=0.6cm, nosep, itemsep=2pt]
    \item Addition
    \item Subtraction
    \item Multiplication
    \item Division
    \item Exponentiation
  \end{enumerate}

\end{enumerate}
\end{multicols}

\vspace{-0.2cm}
\hrule
\vspace{0.2cm}
\textbf{\small Open Ended Questions (FRQ)}
\begin{enumerate}[label=\textbf{Q\arabic*.}, start=9, itemsep=5pt, topsep=0pt]

  \item \small Draw a picture that represents the expression $x + 2$. Be sure to include labels and explain the meaning of each part. \vspace{2cm}

  \item \small Explain the difference between a constant and a variable in your own words, using examples to support your explanation. \vspace{2cm}

\end{enumerate}

\vfill

% Chef's Tips Footer
\begin{tcolorbox}[colback=blue!5, colframe=mainblue!50, title=\small \textbf{Chef's Tips}, fonttitle=\bfseries, bottom=1mm]
\begin{itemize}[noitemsep, topsep=2pt, leftmargin=0.5cm]
  \item \scriptsize Focus on basic conceptual definitions and single-step identification to ensure students understand the fundamentals of variables and expressions.\item \scriptsize Emphasize the importance of understanding the order of operations to simplify expressions correctly.\item \scriptsize Use visual aids and real-world examples to help students connect abstract concepts to tangible scenarios, enhancing their comprehension of algebraic foundations.
\end{itemize}
\end{tcolorbox}

\end{document}
  